\section{模式二：单步评测 --- 高效的``单元测试''}

\subsection{单步评测的核心思想}

LangChain 团队报告，在他们的 Deep Agent 评测用例中，\textbf{约一半}采用单步评测模式。单步评测的核心问题是：

\begin{quote}
\textit{给定特定的输入消息序列，LLM 做出的\textbf{下一步决策}是否正确？}
\end{quote}

这类似于软件工程中的\textbf{单元测试}——不关心整个系统的端到端行为，只验证单个决策点的正确性。

\subsection{适用场景}

单步评测特别适合验证以下类型的决策：

\begin{itemize}
  \item Agent 是否选择了正确的工具来搜索会议时间？
  \item Agent 是否检查了正确的目录内容？
  \item Agent 是否更新了其记忆？
  \item Agent 是否正确地拒绝了不合理的请求？
\end{itemize}

\begin{knowledgebox}{为什么回归问题多发生在单个决策点？}
LangChain 团队的经验表明，Agent 的回归（regression）更多发生在单个决策点而非完整执行序列上。这与软件工程中 bug 倾向于出现在``接口边界''的规律一致。当 prompt 或模型版本更新时，Agent 在特定场景下的工具选择和参数生成最容易出现偏差。
\end{knowledgebox}

\subsection{LangGraph 中断机制}

如果使用 LangGraph 框架，可以利用其 streaming 能力在 Agent 执行第一个工具调用后立即中断，从而高效地实现单步评测：

\begin{lstlisting}[caption={LangGraph 单步评测 --- interrupt\_before 机制}, label={lst:single-step}]
@pytest.mark.langsmith
def test_single_step() -> None:
    state_before_tool_execution = await agent.ainvoke(
        inputs,
        # interrupt_before: 指定在哪些节点前中断
        # 在 tools 节点前中断可以检查工具调用参数
        interrupt_before=["tools"]
    )
    # 获取 Agent 的消息历史，包括最新的工具调用
    print(state_before_tool_execution["messages"])
\end{lstlisting}

\begin{importantbox}{单步评测的效率优势}
单步评测不仅执行速度快（通常只需一次 LLM 调用），还能\textbf{大幅节省 token 消耗}。对于需要频繁运行的回归测试套件来说，这意味着显著更低的成本和更快的反馈循环。LangChain 建议将单步评测作为 CI/CD pipeline 中的第一道关卡。
\end{importantbox}

\subsection{本章小结}

单步评测是 Deep Agent 评测中最高效的手段。它聚焦于单个决策点的正确性，类似于单元测试。通过 LangGraph 的 \texttt{interrupt\_before} 机制可以优雅地实现。建议将约一半的测试用例设计为单步评测，覆盖 Agent 的关键决策分支。


